\begin{pythoncode}
def h_jac(self, eta: np.ndarray) -> np.ndarray:
    """Calculate the jacobian of h.

    Parameters
    ----------
    eta : np.ndarray, shape=(3 + 2 * #landmarks,)
    The robot state and landmarks stacked.

    Returns
    -------
    np.ndarray, shape=(2 * #landmarks, 3 + 2 * #landmarks)
    the jacobian of h wrt. eta.
    """
    # extract states and map
    x = eta[0:3]
    # reshape map (2, #landmarks), m[j] is the jth landmark
    m = eta[3:].reshape((-1, 2)).T

    numM = m.shape[1]

    Rot = rotmat2d(x[2])

    delta_m = m - x[0:2, None]

    zc = delta_m - Rot @ self.sensor_offset[:, None]

    zpred = self.h(eta).reshape(-1, 2).T
    zr = zpred[0]

    Rpihalf = rotmat2d(np.pi / 2)

    # Allocate H and set submatrices as memory views into H
    H = np.zeros((2 * numM, 3 + 2 * numM))
    Hx = H[:, :3]  # slice view, setting elements of Hx will set H as well
    Hm = H[:, 3:]  # slice view, setting elements of Hm will set H as well

    jac_Zcb = -np.eye(2, 3)  # preallocate and update this for speed
    for i in range(numM):  # this whole loop can be vectorized
        ind = 2 * i
        inds = slice(ind, ind + 2)
        jac_Zcb[:, 2] = -Rpihalf @ delta_m[:, i]

        Hx[ind] = (zc[:, i] / zr[i]) @ jac_Zcb
        Hx[ind + 1] = (zc[:, i] / zr[i] ** 2) @ Rpihalf.T @ jac_Zcb

        Hm[inds, inds] = -Hx[inds, :2]

    # H = np.hstack([Hx, Hm]), alternative to using the view

    # TODO: You can set some assertions here 
    # to make sure that some of the structure in H is correct
    return H
\end{pythoncode}